\chapter{Compact coefficients and automorphic products}
\label{mht:ch:compact-examples}

A protected trace retains signed multiplicities of states.
An extension product specifies how states combine.
Geometric field operations also specify the support of their insertions.
For a K3 surface times an elliptic curve, the Igusa product gives an
exact scalar comparison. The constructions below exhibit the additional
algebraic and geometric information needed to compare operators.

We begin with the Jacobi coefficients and their normalization.
They determine a root determinant and a convergent diagonal operator.
We then construct correspondence actions on cohomology and examine
their integral specialization. The Hall correspondence uses a common
elliptic torsor, while its coefficients require a compatible orientation.
Each construction has its own state space and unit.

\section{The Jacobi seed and its scalar normalization}

Write \(q=e^{2\pi i\tau}\) and \(r=e^{2\pi iz}\).
For \(\operatorname{Im}\tau>0\), use the theta series
\[
\begin{aligned}
 \vartheta_2(\tau,z)&=\sum_{n\in\mathbb Z}
             q^{(n+1/2)^2/2}r^{n+1/2},\\
 \vartheta_3(\tau,z)&=\sum_{n\in\mathbb Z}q^{n^2/2}r^n,\\
 \vartheta_4(\tau,z)&=\sum_{n\in\mathbb Z}(-1)^nq^{n^2/2}r^n.
\end{aligned}
\]
The normalized weak Jacobi form of weight zero and index one is
\begin{equation}\label{mht:eq:compact-jacobi}
 \phi_{0,1}(\tau,z)
   =4\sum_{a=2}^4
       \frac{\vartheta_a(\tau,z)^2}{\vartheta_a(\tau,0)^2}
   =\sum_{n\geq0,\,\ell\in\mathbb Z}f(n,\ell)q^nr^\ell.
\end{equation}
The K3 elliptic genus is \(2\phi_{0,1}\).
The coefficients are integral, satisfy \(f(n,\ell)=f(n,-\ell)\),
and vanish for \(4n-\ell^2<-1\).
These are the usual index-one Jacobi conventions.

\begin{proposition}\label{mht:prop:compact-jacobi-coefficients}
For \(0\leq n\leq3\), the coefficients with \(\ell\geq0\) are
\[
 \begin{array}{c|rrrr}
 n\backslash\ell&0&1&2&3\\ \hline
 0&10&1&0&0\\
 1&108&-64&10&0\\
 2&808&-513&108&1\\
 3&4016&-2752&808&-64
 \end{array}.
\]
The negative-\(\ell\) coefficients follow by symmetry, and all
other coefficients in these four rows vanish.
\end{proposition}
\begin{proof}
Put \(t=q^{1/8}\) and \(x=r^{1/2}\).
The three theta series become
\[
 \sum_n t^{4n^2+4n+1}x^{2n+1},\qquad
 \sum_n t^{4n^2}x^{2n},\qquad
 \sum_n(-1)^nt^{4n^2}x^{2n}.
\]
Remove the common factor \(t\) from the numerator and denominator
of the first quotient. Each denominator now has a nonzero constant
coefficient. If \(A(t)=\sum A_jt^j\), \(B(t)=\sum B_jt^j\)
and \(B_0\ne0\), the quotient coefficients are determined by
\[
 C_0=A_0/B_0,\qquad
 C_j=B_0^{-1}\left(A_j-\sum_{i=1}^jB_iC_{j-i}\right).
\]

Squaring the displayed series and applying this recurrence through
\(t^{24}\) gives the table. Terms with exponent not divisible by
eight cancel in the sum. Every surviving \(x\)-exponent is even,
so replacing \(t^8\) by \(q\) and \(x^2\) by \(r\) gives the stated
coefficients.
\end{proof}

For \(Z=\left(\begin{smallmatrix}\tau&z\\z&\sigma\end{smallmatrix}\right)\)
in the Siegel upper half-space, put \(s=e^{2\pi i\sigma}\).
Call a triple \((n,\ell,m)\) positive when \(n,m\geq0\),
either \(n+m>0\), or \(n=m=0\) and \(\ell<0\).
The monic product is
\begin{equation}\label{mht:eq:compact-igusa-product}
 D_5(Z)=q^{1/2}r^{1/2}s^{1/2}
       \prod_{(n,\ell,m)>0}
       (1-q^nr^\ell s^m)^{f(nm,\ell)}.
\end{equation}

The Borcherds--Gritsenko--Nikulin theorem identifies this product
with the weight-five Igusa form, with a character of order two.
Its square is the monic weight-ten form \(\chi_{10}\).
The product converges in a cusp chamber and extends holomorphically
to the Siegel upper half-space.\footnote{
V.~Gritsenko and V.~Nikulin,
\emph{Siegel automorphic form corrections of some Lorentzian
Kac--Moody Lie algebras}.
See \href{https://arxiv.org/abs/alg-geom/9504006}{Theorem~4.1}.

The same authors' paper
\emph{The Igusa modular forms and the simplest Lorentzian
Kac--Moody algebras} gives
\href{https://arxiv.org/abs/alg-geom/9603010}{equations~(1.6)--(1.10)}.
The former paper calls the unscaled theta product \(\Delta_5\).
The latter uses its monic normalization.}

The characteristic phase fixes the sign of the theta product.
For \(a,b\in\{0,1\}^2\) with \(a\cdot b\) even, define
\[
 \Theta_{a,b}(Z)=
 \sum_{\nu\in\mathbb Z^2}
 \exp\!\left(\pi i\left[
  (\nu+a/2)^{\mathsf T}Z(\nu+a/2)+b^{\mathsf T}\nu\right]\right),
 \qquad D_\theta=\prod_{a\cdot b\ {\mathrm{even}}}\Theta_{a,b}.
\]
There are ten factors. This convention uses \(b^{\mathsf T}\nu\).
Replacing it by \(b^{\mathsf T}(\nu+a/2)\) changes the sign of
the complete product.

\begin{proposition}\label{mht:prop:compact-theta-normalization}
With these phases,
\[
 D_\theta=64D_5,\qquad \chi_{10}=2^{-12}D_\theta^2.
\]
The coefficient of \(q^{1/2}s^{1/2}\) in \(D_\theta\) is
\(64(r^{1/2}-r^{-1/2})\).
\end{proposition}
\begin{proof}
For \(a=(0,0)\), all four factors have leading term one.
For \(a=(1,0)\), the two factors have leading term \(2q^{1/8}\).
For \(a=(0,1)\), they have leading term \(2s^{1/8}\).
The two remaining leading expressions are
\[
 2q^{1/8}s^{1/8}(r^{1/4}+r^{-1/4}),\qquad
 2q^{1/8}s^{1/8}(r^{1/4}-r^{-1/4}).
\]
Their product, together with the other factors, is
\(64q^{1/2}s^{1/2}(r^{1/2}-r^{-1/2})\).

The theta-product identity in the cited normalization gives
\(D_\theta=64D_5\), and squaring gives the second equality.
Under the alternate phase convention only the factor
\(a=b=(1,1)\) changes sign.
\end{proof}

The factors with \(m=0\) are determined by \(f(0,0)=10\) and
\(f(0,\pm1)=1\). Thus the first Fourier--Jacobi coefficient of
\(D_5\), after removing \(s^{1/2}\), is
\begin{equation}\label{mht:eq:compact-boundary-factor}
 b(q,r)=q^{1/2}r^{1/2}(1-r^{-1})
    \prod_{n\geq1}(1-q^nr^{-1})(1-q^n)^{10}(1-q^nr).
\end{equation}
The triple product gives \(b=\eta^9\vartheta_{11}\) when
\[
 \vartheta_{11}(\tau,z)=
 \sum_{n\in\mathbb Z}(-1)^n
    q^{(n+1/2)^2/2}r^{n+1/2}.
\]
Its leading factor \(r^{1/2}-r^{-1/2}\) is essential to the
diagonal zero. The divisor of \(D_5\) on the upper half-space
is the orbit of \(z=0\), with multiplicity one.

As an automorphic section it lies in
\(\mathcal L^5\otimes\nu_5\), where \(\mathcal L\) is the Hodge
line and \(\nu_5\) is the order-two character.
The character records the flat factor of this line.

\section{Root degrees, parity and the completed bar}

Let \(L\) have basis \(f_2,f_3,f_{-2}\), with
\((f_2,f_{-2})=-1\), \((f_3,f_3)=2\), and other pairings zero.
The sublattice
\[
 L_{\mathrm{II}}=2\mathbb Z f_2\oplus\mathbb Z f_3
                  \oplus2\mathbb Z f_{-2}
\]
has signature \((2,1)\). Put
\[
 \alpha(n,\ell,m)=2nf_2-\ell f_3+2mf_{-2},\quad
 \delta_1=2f_2-f_3,\quad
 \delta_2=2f_{-2}-f_3,\quad
 \delta_3=f_3.
\]
Direct substitution gives
\[
 (\alpha,\alpha)=-2(4nm-\ell^2),\qquad
 ((\delta_i,\delta_j))=
 \begin{pmatrix}2&-2&-2\\-2&2&-2\\-2&-2&2\end{pmatrix}.
\]
The Weyl vector \(\rho=f_2-\tfrac12f_3+f_{-2}\) satisfies
\((\rho,\delta_i)=-1\).

In the ordered basis \(f_2,f_3,f_{-2}\), the reflections
\(x\mapsto x-(x,\delta_i)\delta_i\) are
\[
 \begin{gathered}
 S_1=\begin{pmatrix}1&4&4\\0&-1&-2\\0&0&1\end{pmatrix},
 \qquad
 S_2=\begin{pmatrix}1&0&0\\-2&-1&0\\4&4&1\end{pmatrix},\\
 S_3=\begin{pmatrix}1&0&0\\0&-1&0\\0&0&1\end{pmatrix}.
 \end{gathered}
\]
Each squares to the identity and preserves the displayed form.
It also preserves \(L_{\mathrm{II}}\).

The Gritsenko--Nikulin construction supplies a Lie superalgebra
\(\mathfrak g\) with these real simple roots and a specified
multiset of imaginary simple roots.
Its generators, parity and relations supply the bracket.
On its positive root spaces the denominator theorem gives
\[
 d_\alpha:=\operatorname{sdim}\mathfrak g_\alpha=f(nm,\ell).
\]
The scalar \(d_\alpha\) need not determine the two parity dimensions.
In particular, adjoining an even and an odd vector of the same
degree leaves every signed multiplicity unchanged.

Write \(\mathfrak n=\bigoplus_{\alpha>0}\mathfrak g_\alpha\).
The positive cone is generated by \(\delta_1,\delta_2,\delta_3\).
In its coordinates
\[
 \alpha=n\delta_1+m\delta_2+(n+m-\ell)\delta_3,\qquad
 h(\alpha)=2n+2m-\ell.
\]
This is a positive integral height on nonzero positive degrees.
Every bounded-height region has finitely many degrees.
Root spaces are finite-dimensional.

For a positive graded object \(V=\bigoplus_{H\geq0}V_H\), its
total-height completion means
\(\widehat V_h=\prod_{H\geq0}V_H\).
Products and coproducts use the completed tensor product graded
by the sum of heights. Every coefficient then involves finitely
many height decompositions.

\begin{proposition}\label{mht:prop:compact-bar}
Super antisymmetrization gives a quasi-isomorphism
\[
 \widehat C_*(\mathfrak n)_h
     \longrightarrow
 \widehat{\operatorname{Bar}(U(\mathfrak n))}_h
\]
of complete coaugmented dg coalgebras.
Its Euler supercharacter is
\[
 \chi_{\mathrm{root}}\widehat C_*(\mathfrak n)_h
        =\prod_{\alpha>0}(1-e^{-\alpha})^{d_\alpha}.
\]
\end{proposition}
\begin{proof}
At fixed total height, only finitely many generators and tensor
lengths occur. The PBW filtration is therefore finite.
The associated graded comparison is the Koszul comparison for
the symmetric superalgebra. An even generator contributes the
usual polynomial Koszul complex. An odd generator contributes
the corresponding exterior-algebra resolution.
Tensoring these one-generator comparisons proves the associated
graded quasi-isomorphism.

The finite filtration proves the
comparison at every height, and exactness of products of vector
spaces proves it after completion.
The antisymmetrization map respects deconcatenation and the
Chevalley coproduct, including the super signs.

For the character use homological degree together with internal
parity. An even root vector contributes \(1-e^{-\alpha}\) to
Chevalley chains, and an odd vector contributes its reciprocal.
Multiplication over a homogeneous basis gives the formula.
\end{proof}

There is also a useful inverse system obtained by first forming
\(\mathfrak n^{(H)}=\mathfrak n/\mathfrak n_{>H}\) and then taking
its ordinary enveloping algebra. It retains all powers of every
surviving generator. Its inverse limit is a different completion.
For the one-dimensional even abelian algebra of height one,
\[
 \varprojlim_HU(\mathfrak n^{(H)})=\mathbb C[x],
 \qquad \widehat{U(\mathfrak n)}_h=\mathbb C[[x]].
\]

Total-height truncation of the entire algebra or bar complex gives
the completion used in the proposition.
The ordinary bar--Chevalley maps also give a comparison of the
generator-quotient inverse systems. That comparison does not
identify the two topologies.

If \(P=\prod_{\alpha>0}(1-e^{-\alpha})^{d_\alpha}\), super PBW gives
\[
 \operatorname{sch}\operatorname{Sym}_{\mathrm{super}}(\mathfrak n)=P^{-1},
 \qquad e^{-\rho}P=\operatorname{Den}_{\mathfrak g}.
\]
In root-torus coordinates \(Q,R,S\), normalize
\(e^{-\alpha(n,\ell,m)}=Q^nR^\ell S^m\).
The chosen Gritsenko--Nikulin embedding into Siegel coordinates
is doubled, so
\[
 \operatorname{Den}_{\mathfrak g}(Z)=D_5(2Z).
\]
The substitution \(Z\mapsto Z/2\) is part of any comparison with
the monic product \(D_5(Z)\).

Elliptic degree \(m\) gives the positive Lie algebra a second
useful decomposition. Its degree-zero part \(\mathfrak n_0\)
is a subalgebra and its positive-degree part \(\mathfrak n_+\)
is an ideal. Root addition proves
\(\mathfrak n=\mathfrak n_0\ltimes\mathfrak n_+\).
PBW, with positive-degree generators ordered first, gives
\[
 \widehat U_h(\mathfrak n)
   \cong\widehat U_h(\mathfrak n_+)\#
                         \widehat U_h(\mathfrak n_0).
\]
Here the smash product uses the adjoint derivation action.

The Chevalley differential splits into the two internal
differentials and that action, giving the total complex
\(\widehat C_*(\mathfrak n_0;C_*(\mathfrak n_+))_h\).
The degree label alone does not turn a point-supported current
into a sheaf whose support covers the elliptic curve.

\section{A convergent root determinant}

For each active degree choose the minimal super vector space
\(V_\alpha\) of signed dimension \(d_\alpha\): it is purely even
if \(d_\alpha\geq0\), and purely odd otherwise.
Give it an orthonormal basis and take the Hilbert direct sum
\(\mathcal H_{\mathrm{vir}}=\bigoplus_\alpha^{\ell^2}V_\alpha\).
This construction uses only the signed multiplicities.
It does not equip \(V_\alpha\) with the root bracket of
\(\mathfrak g\).

\begin{lemma}\label{mht:lem:compact-growth}
There are constants \(C_0,C_1>0\) such that
\[
 |f(N,\ell)|\leq C_0e^{C_1\sqrt{N+1}}.
\]
Consequently \(|d_\alpha|\leq C_2e^{C_3h(\alpha)}\).
\end{lemma}
\begin{proof}
On \(|q|=e^{-t}\), \(|r|=1\), with \(0<t\leq1\), each theta
numerator in \eqref{mht:eq:compact-jacobi} is bounded by
\(Ct^{-1/2}\). For \(\vartheta_2\), first remove its common
factor \(q^{1/8}\). The theta product formulas bound the inverse
denominators by a fixed power of
\(\prod_{j\geq1}(1-e^{-tj/2})^{-1}\), up to a constant.
The estimate
\[
 \sum_{j\geq1}-\log(1-e^{-tj/2})
   =\sum_{k\geq1}\frac1{k(e^{tk/2}-1)}
   \leq\frac2t\sum_{k\geq1}\frac1{k^2}
\]
therefore bounds \(|\phi_{0,1}|\) by \(Ct^{-1}e^{C'/t}\).

Cauchy's coefficient formula in \(q\) and \(r\) gives
\[
 |f(N,\ell)|\leq Ct^{-1}e^{tN+C'/t}.
\]
Set \(t=(N+1)^{-1/2}\) and absorb the polynomial factor.
Finally \(N=nm\), with \(n,m\leq h(\alpha)\), proves the
height bound.
\end{proof}

Let \(\lambda\) be the additive complex energy of a root degree,
so that \(e^{-\alpha(Z)}=e^{-\lambda(\alpha)}\).
Suppose \(\operatorname{Re}\lambda(\delta_i)>0\) for all three
simple roots. If
\(c=\min_i\operatorname{Re}\lambda(\delta_i)\), then
\(\operatorname{Re}\lambda(\alpha)\geq c h(\alpha)\).
At most \(\binom{H+2}{2}\) degrees have height \(H\).
The lemma consequently proves
\[
 \sum_{\alpha>0}|d_\alpha|
       e^{-t\operatorname{Re}\lambda(\alpha)}<\infty
 \quad\hbox{for sufficiently large }t.
\]
Thus \(tZ\) lies in a trace-class region for sufficiently large \(t\).

For the full root superspace the same assertion requires an
estimate on the sum of its two parity dimensions.

In this region define
\[
 K_Z|_{V_\alpha}=e^{-\alpha(Z)}\,\mathrm{id},\qquad D_Z=1-K_Z.
\]
Both parity restrictions of \(K_Z\) are trace class.
The odd determinant is nonzero when all root eigenvalues have
absolute value less than one. Thus
\[
 \operatorname{sdet}_F(D_Z)
    =\frac{\det_F(D_Z|_{\bar0})}{\det_F(D_Z|_{\bar1})}
    =\prod_{\alpha>0}(1-e^{-\alpha(Z)})^{d_\alpha}.
\]
Local uniform summability gives holomorphic dependence in the
trace-class region.

The Pfaffian normalization can be made equally explicit.
On a finite parity block of dimension \(N\), put
\[
 A(D)=\begin{pmatrix}0&D\\-D^{\mathsf T}&0\end{pmatrix},
 \qquad J_N=A(I).
\]
In block order,
\(\operatorname{Pf}A(D)=(-1)^{N(N-1)/2}\det D\).
The relative Pfaffian
\(\operatorname{Pf}A(D)/\operatorname{Pf}J_N\) is therefore
\(\det D\).

Taking the even ratio divided by the odd ratio,
then the trace-norm limit of height truncations, gives
\[
 \operatorname{sPf}_{\mathrm{rel}}(A(D_Z))
       =\operatorname{sdet}_F(D_Z),\qquad
 e^{-\rho}\operatorname{sPf}_{\mathrm{rel}}(A(D_Z))=D_5(2Z).
\]
This is a normalized Fredholm model on a specified Hilbert
superspace. The operator \(1-K_Z\) is bounded.
A first-order geometric operator requires its own symbol,
domain and comparison of determinant lines.
Continuing the scalar automorphic section supplies its
automorphy law. An operator intertwiner is additional data.

\section{The reduced enumerative scalar}

Let \(X=S\times E\), where \(S\) is a nonsingular projective
K3 surface and \(E\) is an elliptic curve.
Choose a nonzero primitive class \(\beta_h\) with
\(\beta_h^2=2h-2\), and \(\beta_h^\vee\in H^2(S,\mathbb Q)\)
with \(\langle\beta_h,\beta_h^\vee\rangle=1\).
If \(p\in H^2(E,\mathbb Z)\) is the point class, define
\[
 N_{g,h,d}
  =\int_{[\overline{\mathcal M}^{\,\bullet}_{g,1}
              (X,(\beta_h,d))]^{\mathrm{red}}}
       \operatorname{ev}_1^*
       \bigl(\pi_S^*\beta_h^\vee\cup\pi_E^*p\bigr).
\]
The bullet allows disconnected domains, with the standard
noncontracted-component convention. The reduced virtual class
and the incidence insertion are part of this definition.

Oberdieck--Pixton's theorem gives
\[
 \sum_{g,h,d\geq0}N_{g,h,d}\,u^{2g-2}q^{h-1}s^{d-1}
     =-\frac1{\chi_{10}(q,e^{iu},s)}.
\]
This is a Laurent expansion under \(r=e^{iu}\), in the sector
primitive in the K3 factor.\footnote{
G.~Oberdieck and A.~Pixton,
\emph{Holomorphic anomaly equations and the Igusa cusp form
conjecture}. See \href{https://arxiv.org/abs/1706.10100}
{Section~0.2, Theorem~1, and Sections~4.4--4.7}.}
The minus sign and the shifts by \(h-1,d-1\) belong to the
stated enumerative normalization.

Let \(\mathcal F=\operatorname{Sym}_{\mathrm{super}}(\mathfrak n)\).
The preceding root calculation gives the exact scalar identity
\[
 -\left.(e^{2\rho}\operatorname{sch}(\mathcal F^{\otimes2}))
                      \right|_{Z\mapsto Z/2}
       =-\chi_{10}(Z)^{-1}.
\]
It follows from \(\operatorname{sch}\mathcal F=P^{-1}\) and
\(e^{-\rho}P=D_5(2Z)\).
The two sides have independently specified sources.
An identification of their states must construct a map carrying
the geometric operators and their trace to the root representation.

\section{Surface coefficients and their integral specialization}

Let \(Y\) be a projective K3 or Enriques surface.
Use rational cohomology and its complex orientation.
Then \(H^{\mathrm{odd}}(Y)=0\), \(c_1(T_Y)=0\) rationally, and
\[
 A=H^*(Y,\mathbb Q)
   =\mathbb Q1\oplus H^2(Y,\mathbb Q)\oplus\mathbb Q\omega,
 \qquad \int_Y\omega=1.
\]
If \(b=b_2(Y)\), then \(c_2(T_Y)=e\omega\), with \(e=b+2\).
Thus \((b,e)=(22,24)\) for K3 and \((10,12)\) for Enriques.
The parameter \(y\) below records the normalized Hodge charge,
and \(q\) records elliptic energy.

The elliptic characteristic class is the Todd class times the
Chern character of
\[
 \mathcal E_{q,y}(Y)=
 y^{-1}\Lambda_{-y}T_Y^*
 \bigotimes_{n\geq1}
 \left(
 \operatorname{Sym}_{q^n}(T_Y\oplus T_Y^*)
 \otimes\Lambda_{-yq^n}T_Y^*
 \otimes\Lambda_{-y^{-1}q^n}T_Y
 \right).
\]
Here \(\operatorname{Sym}_t\) and \(\Lambda_t\) denote the
generating series of symmetric and exterior powers.
Write \(\epsilon_Y=\operatorname{td}(T_Y)
\operatorname{ch}(\mathcal E_{q,y}(Y))\).
Its integral is the elliptic genus.
This class carries more information than its integral.

\begin{proposition}\label{mht:prop:compact-surface-class}
Put \(s_y=y+y^{-1}-2\) and \(a_y=1+s_y/12\).
Then
\[
 [q^0]\epsilon_Y=s_y\,1+a_yc_2(T_Y),\qquad
 [q^1]\epsilon_Y=-2s_y^2\,1+
                 \left(\frac56s_y^2-2s_y\right)c_2(T_Y).
\]
For Enriques, integration of the second expression is
\(10y^2-64y+108-64y^{-1}+10y^{-2}\).
\end{proposition}
\begin{proof}
Use the formal opposite Chern roots \(x,-x\), so \(c_2=-x^2\).
Modulo classes above degree four, the Todd factor is
\(1-x^2/12\), and
\[
 y^{-1}(1-ye^x)(1-ye^{-x})=s_y-x^2.
\]
This proves the constant coefficient.

The first oscillator coefficient is
\[
 (2-y-y^{-1})(e^x+e^{-x})=-2s_y+s_yc_2.
\]
Multiplying it by the constant coefficient gives the
\(q\)-linear formula. Integrating with \(\int_Yc_2=12\)
gives the displayed Laurent polynomial.
For K3 the integral is twice that polynomial.
\end{proof}

Now specialize to \(q=0\), put \(t=y-1\), and use
\(R_0=\mathbb Q[[t]]\).
Multiplication by
\(\epsilon=s_y+a_yc_2\) acts on
\(H=A\otimes R_0\), where
\(s_y=t^2/(1+t)\).
It is invertible over \(\mathbb Q((t))\) and injective over
\(R_0\). Its image \(\mathcal M=\epsilon H\) is a free lattice.

\begin{proposition}\label{mht:prop:compact-smith}
The cokernel of this multiplication has the decomposition
\[
 \operatorname{coker}(\epsilon:H\to H)
   \cong R_0/(t^4)\oplus\bigl(R_0/(t^2)\bigr)^b.
\]
Its derived fiber at \(t=0\) has dimension \(b+1\) in
each of degrees \(-1,0\), and vanishes in other degrees.
The image of \(\mathcal M\) in \(H/tH\) has rank one.
\end{proposition}
\begin{proof}
On \(H^2(Y)\), multiplication is the scalar \(s_y\), a unit
times \(t^2\). On the ordered basis \(1,\omega\), the matrix is
\[
 M=\begin{pmatrix}s_y&0\\u&s_y\end{pmatrix},\qquad u=a_ye.
\]
The coefficient \(u\) is a unit in \(R_0\).
Swap the rows, subtract \(s_y/u\) times the first row from
the second, and then subtract \(s_y/u\) times the first
column from the second. The resulting matrix is
\(\operatorname{diag}(u,-s_y^2/u)\).
Its Smith form is \(\operatorname{diag}(1,t^4)\).

Each module \(R_0/(t^k)\) has free resolution
\([R_0\xrightarrow{t^k}R_0]\).
Derived tensor with \(R_0/(t)\) makes its differential zero.
There are \(b+1\) torsion summands.
Finally \(\epsilon(0)=e\omega\), whose multiplication has
one-dimensional image.
\end{proof}

The free module \(\mathcal M/t\mathcal M\) has dimension \(b+2\).
Its map into \(H/tH\) need not be injective.
The exact sequence is
\[
 0\longrightarrow\operatorname{Tor}_1^{R_0}
        (\operatorname{coker}\epsilon,\mathbb Q)
 \longrightarrow\mathcal M/t\mathcal M
 \longrightarrow H/tH
 \longrightarrow(\operatorname{coker}\epsilon)/t
 \longrightarrow0.
\]
The inverse outer block also detects a pole:
\[
 M^{-1}=
 \begin{pmatrix}s_y^{-1}&0\\-a_ye\,s_y^{-2}&s_y^{-1}\end{pmatrix},
 \qquad
 \int_Y\epsilon^{-1}=-a_ye\,s_y^{-2}.
\]
This is a fourth-order pole at \(t=0\), although
\(\int_Y\epsilon=a_ye\) is regular.

\section{Correspondence actions and the elliptic state}

Put \(X_n=Y^{[n]}\), with \(X_0=\mathrm{pt}\).
A cohomological kernel from \(X_m\) to \(X_n\) is
\(K\in H^*(X_n\times X_m,\mathbb Q)\).
It acts by
\[
 T_K(v)=(p_n)_*(K\cup p_m^*v).
\]
All spaces are smooth and projective, so these pushforwards
are proper. K\"unneth and Poincar\'e duality identify the kernel
space, with its source-dimension shift, with all graded linear
maps between the two cohomologies.

Pullback, cup product and integration over the middle factor
therefore give associative convolution and diagonal identities.
The representation-theoretic information lies in the selected
geometric kernels.

For \(r>0\), the incidence correspondence records
\[
 Z\subset Z',\quad \ell(Z)=n,\quad \ell(Z')=n+r,\quad
 \operatorname{Supp}(I_Z/I_{Z'})=\{x\}.
\]
Its fundamental cycle, with a class \(a\in A\) inserted at
\(x\), defines creation \(C_r(a)\).
For the ordinary Poincar\'e pairing define
\(A_r(a)=(-1)^{r-1}C_r(a)^\dagger\).

The geometric Heisenberg theorem and Fock isomorphism give
\[
 [A_r(a),C_s(b)]=r\delta_{rs}\int_Yab,\qquad
 \operatorname{Sym}\!\left(\bigoplus_{r\geq1}p^rA\right)
     \cong\bigoplus_{n\geq0}H^*(X_n,\mathbb Q).
\]
Creators of the same sign commute, as do annihilators.

The Fock map sends a polynomial in \(p_r(a)\) to the same
polynomial in \(C_r(a)\) applied to the empty vacuum.\footnote{
M.~Lehn and C.~Sorger,
\emph{The cup product of the Hilbert scheme for K3 surfaces}.
See \href{https://arxiv.org/abs/math/0012166}
{Theorems~3.2, 3.5--3.6 and Proposition~3.9}.
Their surface trace is \(-\int_Y\).
The annihilators here use \(+r\int_Yab\).}

The normalization of the degree-zero unit is
\[
 \frac1{n!}C_1(1)^n|0\rangle=1_{X_n}.
\]
Over \(n\) distinct points the iterated incidence orders them
in \(n!\) ways. Both sides have degree zero, which proves
the equality.

Let \(\epsilon_n=\operatorname{td}(T_{X_n})
\operatorname{ch}(\mathcal E_{q,y}(X_n))\), with
the factor \(y^{-n}\) appropriate to complex dimension \(2n\).
Use \(R=\mathbb Q(y)[[q]]\).
Its degree-zero component \(r_n\) has constant coefficient
\(y^{-n}(1-y)^{2n}\), a unit in \(R\).
The positive-degree part \(\nu_n=\epsilon_n-r_n\) satisfies
\(\nu_n^{2n+1}=0\), so
\[
 \epsilon_n^{-1}
    =r_n^{-1}\sum_{j=0}^{2n}(-r_n^{-1}\nu_n)^j.
\]
This is an inverse in the cohomology ring.
Multiplication by it need not preserve each ordinary
cohomological degree.

\begin{proposition}\label{mht:prop:compact-transport}
For every cohomological kernel define
\[
 \widetilde K=p_n^*\epsilon_n\cup K\cup p_m^*\epsilon_m^{-1}.
\]
If \(U_j\) is multiplication by \(\epsilon_j\), then
\[
 T_{\widetilde K}=U_nT_KU_m^{-1},\qquad
 \widetilde{L\circ K}=\widetilde L\circ\widetilde K,\qquad
 \widetilde{[\Delta_{X_n}]}=[\Delta_{X_n}].
\]
These identities hold for every finite composable string.
\end{proposition}
\begin{proof}
In the convolution, the two characteristic classes on the
middle factor multiply to one.
The projection formula carries the two remaining factors
through the pushforward. This proves the composite and
operator identities. Pullback to the diagonal proves the
identity-kernel formula. Repeating the same cancellation
proves the assertion for longer strings.
\end{proof}

Complete the state space in length and let
\(\langle\mathcal I|\) integrate coefficientwise in the length variable \(p\).
Since \(U_0=1\), the incidence normalization gives
\[
 \frac1{n!}\widetilde C_1(1)^n|0\rangle=\epsilon_n,\qquad
 \left\langle\mathcal I\middle|
     \exp(p\widetilde C_1(1))\middle|0\right\rangle
        =\sum_{n\geq0}p^n\operatorname{Ell}(X_n;q,y).
\]
The scalar is a specified matrix coefficient of a prepared
state. Conjugation itself leaves matrix traces unchanged.

More generally, any even invertible classes \(\alpha_n\)
with \(\alpha_0=1\) give the same construction with
\(\int_{X_n}\alpha_n\) as coefficients.
For \(n>0\), the choice
\(\alpha_n=1+c_n[\mathrm{pt}]\) realizes any prescribed
coefficient \(c_n\). Thus the scalar alone does not identify
the geometric action.

At \(q=0\), set
\(\mathcal M_n=\epsilon_nH^*(X_n,\mathbb Q)[[t]]\).
For a kernel over \(\mathbb Q((t))\), the transported
operator preserves these image lattices exactly when
the original operator preserves the original
\(\mathbb Q[[t]]\)-modules. Indeed,
\[
 \widetilde T_K(\epsilon_m v)=\epsilon_nT_K(v).
\]
Kernels defined over \(\mathbb Q\) satisfy this condition.
The operator \(t^{-1}\mathrm{id}\) does not.
Proposition~\ref{mht:prop:compact-smith} describes the
nontrivial specialization even in length one.

\section{Joining cycles and the diagonal class}

Choose a basis \(u_i\) of \(H^2(Y,\mathbb Q)\), with
\(Q_{ij}=\int_Yu_iu_j\) and inverse matrix \(Q^{ij}\).
The product satisfies \(u_iu_j=Q_{ij}\omega\).
Let \(\Delta:A\to A\otimes A\) be adjoint to multiplication
for the trace \(\int_Y\). It is the diagonal Gysin map.
Testing against basis pairs gives
\[
\begin{aligned}
 \Delta1&=1\otimes\omega+\omega\otimes1+
                         \sum_{i,j}Q^{ij}u_i\otimes u_j,\\
 \Delta u_i&=u_i\otimes\omega+\omega\otimes u_i,\qquad
 \Delta\omega=\omega\otimes\omega.
\end{aligned}
\]

Consequently \(m\Delta1=(b+2)\omega=e_A=c_2(T_Y)\)
and \(m\Delta(a)=e_Aa\).
A split followed by a join inserts this Euler class.

For permutations \(\sigma,\tau\in\mathfrak S_N\), label
each cycle by a class in \(A\).
On a joint orbit \(O\) of their action set
\[
 g_O=\frac{|O|+2-c_\sigma(O)-c_\tau(O)
                              -c_{\sigma\tau}(O)}2.
\]
The transitive triple
\(\sigma,\tau,(\sigma\tau)^{-1}\) defines a branched cover
of the sphere, so Riemann--Hurwitz makes \(g_O\) a
nonnegative integer.
Multiplying the input labels to obtain \(a_O,b_O\), the
Frobenius permutation product on that orbit is
\[
 \Delta_{c_{\sigma\tau}(O)}(a_Ob_Oe_A^{g_O}),
\]
where \(\Delta_k\) is the iterated diagonal and
\(\Delta_1=\mathrm{id}\).
The projection formula and \(m\Delta(a)=e_Aa\) make this
operation independent of a decomposition of the covering
surface. They give the associativity of the Frobenius
permutation algebra.

For \(s=(12)\) and \(\rho=(123)\), the first products are
\[
 (as)(bs)=\Delta(ab)\,\mathrm{id},\qquad
 (a\rho)(b\rho)=abe_A\rho^2,\qquad
 (a\rho)(b\rho^{-1})=\Delta_3(ab)\,\mathrm{id}.
\]
The middle case has genus one. A joint orbit of genus at
least two gives zero because \(e_A^2=0\).

For the Hilbert-scheme product use the additional sign
\[
 \varepsilon(\sigma,\tau)=
 (-1)^{(|\sigma|+|\tau|-|\sigma\tau|)/2},\qquad
 |\sigma|=N-c(\sigma).
\]
Permutation parity makes the exponent integral.
Its telescoping identity is the two-cocycle equation.
The sign therefore preserves associativity.

The Lehn--Sorger theorem identifies the resulting invariant
algebra with \(H^*(Y^{[N]},\mathbb Q)\) when the canonical
class is numerically trivial, including both surfaces used
here. The sector factor \(i^{|\sigma|}\) removes the sign
over \(\mathbb C\). Over \(\mathbb Q\) the sign remains explicit.

This product has a useful differential-operator expression.
In a homogeneous basis \(a_\alpha\) write
\[
 a_\alpha a_\beta=m_{\alpha\beta}^{\gamma}a_\gamma,\qquad
 \Delta a_\gamma=\Delta_\gamma^{\alpha\beta}
                              a_\alpha\otimes a_\beta.
\]
On the Fock polynomial algebra let \(\partial_{n,\alpha}\)
differentiate \(p_n(a_\alpha)\), and define
\begin{equation}\label{mht:eq:compact-join}
\begin{aligned}
 J&=\frac12\sum_{i,j\geq1}ij\,m_{\alpha\beta}^{\gamma}
              p_{i+j}(a_\gamma)\partial_{i,\alpha}\partial_{j,\beta},\\
 S&=\frac12\sum_{i,j\geq1}(i+j)\Delta_\gamma^{\alpha\beta}
              p_i(a_\alpha)p_j(a_\beta)\partial_{i+j,\gamma},
 \qquad D_H=J-S.
\end{aligned}
\end{equation}

Every fixed-length component has finitely many terms.
There are \(ij\) choices of endpoints joining two cycles.
There are \(i+j\) positions for a split, with the ordered
sum counting each split twice.
The Hilbert sign gives \(D_H\), which is multiplication by
\(-c_1(\mathcal O_Y^{[N]})\) under the Fock isomorphism.

\begin{proposition}\label{mht:prop:compact-joining}
For multiplication \(M_i(a)\) by \(p_i(a)\),
\[
 [[D_H,M_i(a)],M_j(b)]=ijM_{i+j}(ab),\qquad
 [[[D_H,M_i(a)],M_j(b)],M_k(c)]=0.
\]
Moreover
\[
 D_H^2p_2(a)=-p_2(e_Aa),\qquad
 D_H^{2N+1}|_{H^*(X_N)}=0.
\]
\end{proposition}
\begin{proof}
The splitting term is first order, so its double commutator
with multiplication vanishes.
The two derivative slots in \(J\) give the first formula.
Its right side is multiplication, giving the third
commutator identity.

At length two, splitting \(p_2(a)\) and then joining gives
\(-p_2(m\Delta a)\).
Both operations raise ordinary cohomological degree by two.
A join raises the cycle-age shift by two.
A split inserts a degree-four diagonal and lowers that shift by two.
The cohomology of \(X_N\) lies between degrees zero and
\(4N\), proving the nilpotence bound.
\end{proof}

In particular \(D_H^2p_2(1)\ne0\).
Joining is an even operator of order two.
Its double commutator measures its failure to be a derivation.
These are cohomological diagonal operations.
Lifting them to a chiral state complex requires supported
maps retaining the normal oscillators and their differential.

\section{Extension products and their coefficients}

An elementary Hall action already exists on cyclic quotients
of a surface. Let \(\mathcal C_0(Y)\) be the category of
zero-dimensional coherent sheaves on \(Y\).
For constructible functions, with Euler integration over
subobjects of fixed length, set
\[
 (h_1*h_2)(Q)=\int_{U\subset Q}
                  h_1(U)h_2(Q/U)\,d\chi.
\]
The order here places the subobject first.
A cyclic quotient is a surjection
\(\mathcal O_Y\twoheadrightarrow Q\), and its quotient by
any \(U\subset Q\) remains cyclic.
The formula
\[
 (h\cdot v)(\mathcal O_Y\twoheadrightarrow Q)
   =\int_{U\subset Q}h(U)
                   v(\mathcal O_Y\twoheadrightarrow Q/U)\,d\chi
\]
therefore defines a left action.

The two parenthesizations integrate the same function
\[
 h_1(U)h_2(W/U)
           v(\mathcal O_Y\twoheadrightarrow Q/W)
\]
over \(0\subset U\subset W\subset Q\).
The subobject spaces are projective Quot varieties, and
Euler pushforward composes. This proves associativity,
all finite-flag compatibilities, and the zero-object unit.
Positive-length completion is allowed because each length
has finitely many decompositions.

The construction uses constructible functions.
A critical cohomological Hall action requires its own
pullback, pushforward and orientation maps.
There is also a change of category on passing from \(Y\)
to \(Y\times E\). In a local normal coordinate \(z\), an
extension of two modules killed by \(z\) can have
\[
 z=\begin{pmatrix}0&1\\0&0\end{pmatrix}.
\]
Thus sheaves with zero normal operator need not be closed
under the ambient extensions.

For compact sheaves modulo elliptic translation, even the
base of the extension correspondence matters.
Let \(BE\) be the stack of right \(E\)-torsors on the fppf
site. For a left \(E\)-stack \(M\), an object of \([M/E]\)
is a torsor \(P\) together with an equivariant map \(P\to M\).
The forgetful map lands in \(BE\).

\begin{proposition}\label{mht:prop:compact-common-torsor}
There is a natural equivalence
\[
 [(M\times N)/E_{\mathrm{diag}}]
       \simeq[M/E]\times_{BE}[N/E].
\]
The finite-factor equivalences commute with regrouping.
For coherent sheaves on \(S\times E\), the fixed-polynomial
extension spans over \(BE\) form a coherently associative
unital correspondence object.
The unit is the zero-sheaf map \(BE\to[\mathcal M/E]\).
\end{proposition}
\begin{proof}
An object of the fiber product consists of
\((P,m),(Q,n)\) and an isomorphism \(\alpha:P\to Q\).
It gives the single torsor \(P\) with maps \(m,n\alpha\).
The inverse uses the same torsor twice and the identity.
On morphisms, compatibility with \(\alpha\) determines
one torsor map from the other. This proves the equivalence,
including descent and base change.
Several factors give one torsor with the same list of maps,
which proves the regrouping coherences.

Fix the Hilbert polynomials of the two ends of an exact
sequence. Over a flat family \(F\), its quotients of the
specified polynomial form a projective Quot scheme after
trivializing the torsor. The kernel is flat because both
the middle and quotient families are flat.
An automorphism fixing the middle object and quotient
is the identity. Descent therefore gives a representable
proper map from the extension stack to the middle-object
stack.

Both composites of three extension spans classify
\(0\subset A\subset D\subset F\) on that same torsor.
Recovering \(D\) from a subobject of \(F/A\) uses its
inverse image in \(F\). This is inverse to passage to the
quotient, commutes with base change, and preserves flatness.
Every longer regrouping gives the same flag.
The zero sheaf on a torsor gives the stated unit.
\end{proof}

For \(M=N=E\) with translation, the diagonal quotient is
\(E\), with coordinate \(y-x\), whereas the two separate
quotients are points. The fiber product over \(BE\)
retains their relative translation.
A stability restriction can replace the Quot fiber by an
open subspace, which need not remain proper.
A coefficient theory must also include the nonaffine
stabilizer \(E\), or be defined relative to \(BE\).

The determinant calculation explains the orientation data.
For a three-Calabi--Yau dg category, assume that the relative
Hom complexes \(C(A,B)\) are perfect and that the trace gives
\[
 C(B,A)\simeq C(A,B)^\vee[-3]
\]
compatibly with base change.
Put \(L(A,B)=\det C(A,B)\), \(D(A)=L(A,A)\), with
\(\det C=\bigotimes_i(\det C^i)^{(-1)^i}\).
The symmetry of determinant lines carries parity
\(\chi(C)\bmod2\).

\begin{proposition}\label{mht:prop:compact-extension-line}
For a triangle \(A\to F\to B\to A[1]\),
\[
 D(F)\simeq D(A)\otimes D(B)\otimes L(A,B)^{\otimes2}.
\]
The determinant comparisons agree under refinement of flags,
with the stated symmetry signs.
\end{proposition}
\begin{proof}
Apply derived Hom to the triangle in both variables.
Determinant additivity expresses \(D(F)\) as the tensor
product of the four complexes \(C(A,A),C(A,B),C(B,A),C(B,B)\).
Dualization and the odd shift each invert the determinant
line, so their composite identifies \(L(B,A)\) with
\(L(A,B)\). Reorder with determinant symmetry.
For a longer flag all decompositions give the same ordered
product over pairs of successive quotients.
Compatibility of determinant additivity with refinement
proves the assertion.
\end{proof}

An orientation therefore includes lines \(o_A\) with
\(o_A^{\otimes2}\simeq D(A)\) and maps
\[
 o_A\otimes o_B\otimes L(A,B)\longrightarrow o_F
\]
whose squares are the determinant comparisons.
The two composites on a three-step flag can differ by
an element of \(\mu_2\). Its coherent trivialization is
additional to the square identities.

The cross line can itself be nontrivial.
For \(T=\mathbb P^1\), \(x\in X\), and
\[
 A=\mathcal O_X\boxtimes\mathcal O_T(1),\qquad
 B=\mathcal O_x\boxtimes\mathcal O_T,
\]
one has \(C(A,B)=\mathcal O_T(-1)[0]\) and
\(C(B,A)=\mathcal O_T(1)[-3]\).
Hence the cross factor in \(D(A\oplus B)\) is
\(\mathcal O_T(-2)\), although the self-determinants
are constant lines. Its orientation factor has degree
\(-1\), and a scalar sign cannot replace it.

\section{Physical charge and the Gram map}

The Mukai lattice of a K3 surface is
\[
 \Lambda=H^0(S,\mathbb Z)\oplus H^2(S,\mathbb Z)
                         \oplus H^4(S,\mathbb Z),\qquad
 ((r,c,s),(r',c',s'))=c\cdot c'-rs'-r's.
\]
It is even unimodular of rank \(24\) and signature \((4,20)\).

For electric and magnetic charges \(c=(Q,P)\in\Lambda\oplus\Lambda\),
define
\[
 \Pi(c)=\left(Q^2/2,Q\cdot P,P^2/2\right)\in\mathbb Z^3.
\]
Physical charge adds. The Gram map is quadratic:
\[
\begin{aligned}
 \Pi(c+d)&=\Pi(c)+\Pi(d)+B(c,d),\\
 B((Q,P),(Q',P'))&=
       (Q\cdot Q',\,Q\cdot P'+Q'\cdot P,\,P\cdot P').
\end{aligned}
\]
Bilinearity proves the cocycle equation
\(B(c,d)+B(c+d,e)=B(d,e)+B(c,d+e)\).
Thus
\[
 (c,\lambda)(d,\mu)
     =(c+d,\lambda+\mu+B(c,d))
\]
defines an abelian extension of the charge lattice by
\(\mathbb Z^3\).

The graph map \(j_\Pi(c)=(c,\Pi(c))\) is a homomorphism.
This extension is split: the explicit isomorphism to the
ordinary direct product is
\[
 (c,\lambda)\longmapsto(c,\lambda-\Pi(c)).
\]
It records the quadratic degree without changing the
additive physical charge.

For a charge-graded algebra \(A=\bigoplus_cA_c\), use
central symbols \(u^\lambda\) and the product
\[
 (a_cu^\lambda)\star(b_du^\mu)
       =a_cb_du^{\lambda+\mu+B(c,d)}.
\]
The cocycle equation proves associativity, and
\[
 a_c\longmapsto a_cu^{\Pi(c)}
\]
is an algebra map. The splitting sends this graph copy
back to \(A\otimes1\).
For a completed Hall algebra the charge map and splitting
must also preserve the chosen completion.
This construction regrades a supplied Hall integration
map. It does not construct its coefficient theory.

The common K3 state space also carries several distinct
geometric actions. Nakajima creation acts between Hilbert
schemes of different lengths.
The Lie algebra generated by Lefschetz triples acts on
each fixed hyperk\"ahler cohomology.
Its real form is \(\mathfrak{so}(4,20)\) for K3, and
\(\mathfrak{so}(4,21)\) for \(K3^{[n]}\) type with \(n\geq2\).

Local ADE correspondences use a curve-supported moduli
problem and its equivariant coefficient field.
These statements specify different carriers.
A common quantum algebra requires maps of these carriers,
mixed convolution relations, and a compatible coproduct
and pairing.

For clarity, fix a direct sum \(V\) of the chosen cohomology
spaces and a set of operators preserving finite support.
The free algebra on their symbols maps to
\(\operatorname{End}(V)\). Quotient by its kernel gives the
algebra they generate.
This faithful action does not determine a tensor structure.

If a rigid linear tensor category and a strong monoidal
finite-dimensional fiber functor \(\omega\) are separately
specified, their matrix coefficients instead form the coend
\[
 H=\left(\bigoplus_M\omega(M)^*\otimes\omega(M)\right)
 \Big/
 \left\langle[\lambda\omega(f)\otimes v]_M
       -[\lambda\otimes\omega(f)v]_N\right\rangle.
\]
Its coproduct and counit are
\[
 \Delta[\lambda\otimes v]_M
   =\sum_i[\lambda\otimes e_i]_M\otimes[e^i\otimes v]_M,
 \qquad \epsilon[\lambda\otimes v]_M=\lambda(v).
\]

The identity tensor proves basis independence and
coassociativity. Tensor product supplies multiplication,
and the evaluation and coevaluation identities supply
the antipode. Each selected categorical morphism
intertwines the resulting coactions.
The action algebra records products of selected operators.
The coend records natural symmetries of the specified
fiber functor.

\section{Finite Hall presentations and super relations}

A finite symmetric even Borcherds--Cartan matrix
\(A=(a_{ij})\) determines an ordinary quiver as follows.
Order the vertices. For \(i<j\), place \(-a_{ij}\) arrows
from \(i\) to \(j\). At an imaginary vertex place
\(1-a_{ii}/2\) loops. A real vertex has none.

If
\[
 \langle\mathbf d,\mathbf e\rangle
     =\sum_i d_ie_i-\sum_{a:i\to j}d_ie_j,
\]
then its symmetrization is \(A\).
For the three real roots above the quiver has two arrows
between each ordered pair of distinct vertices.

Over \(\mathbb F_q\), put \(v^2=q\) and use the twisted
subobject-count product
\[
 [M]*[N]=v^{\langle\dim M,\dim N\rangle}
       \sum_{[L]}
       \#\{N'\subset L:N'\simeq N,\ L/N'\simeq M\}[L].
\]
This convention places the quotient first.
Associativity counts the same extension flags in two
orders, and bilinearity makes their twist exponents agree.
The class of zero is the unit.
The Hall \(A_2\) calculation of
Chapter~\ref{ch:algebraic-reconstruction-examples} fixes
the root normalization \(v^{-1}[M]\).

Select the one-dimensional simple at each vertex on which
all loops act by zero.
The generic composition algebra generated by those classes
is the ordinary positive quantum generalized Kac--Moody
algebra of the even matrix.
The Kang--Schiffmann parameter is \(v_{\mathrm{KS}}=v^{-1}\)
for the displayed Hall convention.
Their canonical basis uses the specified semisimple
perverse-sheaf construction on nilpotent representation
varieties.\footnote{
S.-J.~Kang and O.~Schiffmann,
\emph{Canonical bases for quantum generalized Kac--Moody algebras}.
See \href{https://arxiv.org/abs/math/0311089}
{Section~2.3, Theorem~2.2 and Theorem~4.1}, with the
\href{https://arxiv.org/abs/0711.1948}{addendum}.}

The chosen generators are part of this theorem.
A vertex with \(g\) loops has \(q^g\) distinct
one-dimensional simples over \(\mathbb F_q\).
The full equivariant derived category is larger still:
at a real vertex in dimension one its stack is
\(B\mathbb G_m\), with cohomology \(\mathbb Q[c_1]\).
Its total dimension is not the corresponding root multiplicity.

The Borcherds superalgebra has its own labeled simple
generators and parity-dependent relation ideal.
For a single even abelian generator the enveloping algebra
is \(\mathbb C[x]\).
For a single odd abelian generator it is
\(\Lambda(\xi)\), since \(2\xi^2=[\xi,\xi]=0\).
Reassigning parity to an ordinary PBW vector space leaves
its degree-two vector present.
Thus the ordinary Hall presentation and a quantum
supergroup require an explicit comparison respecting
relations and both parity dimensions.

The same quiver has a zero-potential Ginzburg dg algebra.
Use left-to-right path concatenation in this paragraph.
Add reverse arrows \(a^*\) of degree \(-1\) and vertex
loops \(t_i\) of degree \(-2\), with
\[
 da=da^*=0,\qquad
 dt_i=\sum_{s(a)=i}aa^*-\sum_{t(a)=i}a^*a.
\]
The differential squares to zero.
The sum over vertices is a sum of commutators, so it
vanishes in the cyclic quotient and after matrix trace.
It need not vanish in the path algebra.

Projection to a full subquiver kills the omitted vertices
and incident arrows and commutes with this differential.
The reverse inclusion generally fails: adjoining an arrow
changes \(dt_i\) by \(aa^*\) at its old source vertex.
These dg projections are different maps from extension
by zero in the ordinary representation category.

\section{The eight diagonal-divisor forms}

The scalar comparison has a finite arithmetic extension.
For positive integers \(t,N\), define the Hecke subgroup
of the paramodular group by
\[
 \Gamma_t(N)=\operatorname{Sp}_4(\mathbb Q)\cap
 \begin{pmatrix}
 \mathbb Z&t\mathbb Z&\mathbb Z&\mathbb Z\\
 \mathbb Z&\mathbb Z&\mathbb Z&t^{-1}\mathbb Z\\
 N\mathbb Z&Nt\mathbb Z&\mathbb Z&\mathbb Z\\
 Nt\mathbb Z&Nt\mathbb Z&t\mathbb Z&\mathbb Z
 \end{pmatrix}.
\]
Write \(\Gamma_t=\Gamma_t(1)\).
For \(t=1\), this is the usual genus-two
\(\Gamma_0^{(2)}(N)\).

A diagonal-divisor form vanishes simply on the translates
of \(z=0\) and has no other divisor on the upper half-space.
The Cl\'ery--Gritsenko classification gives, up to nonzero
scalar, exactly the following eight forms in this class:
\[
 \begin{array}{c|c|c|c}
 \text{form}&\text{group}&\text{weight}&
          \text{character or multiplier order}\\ \hline
 \Delta_5&\Gamma_1&5&2\\
 \Delta_2&\Gamma_2&2&4\\
 \Delta_1&\Gamma_3&1&6\\
 \Delta_{1/2}&\Gamma_4&1/2&8\\
 \nabla_3&\Gamma_1(2)&3&2\\
 \nabla_2&\Gamma_1(3)&2&2\\
 \nabla_{3/2}&\Gamma_1(4)&3/2&4\\
 Q_1&\Gamma_2(2)&1&4
 \end{array}.
\]
The classification concerns these Hecke subgroups and
integral or half-integral weights.\footnote{
F.~Cl\'ery and V.~Gritsenko,
\emph{The Siegel modular forms of genus 2 with the simplest divisor}.
See \href{https://arxiv.org/abs/0812.3962}
{Theorem~1.4 and Section~3}.}
Half-integral weights use the specified multiplier system.

At higher level the Borcherds product uses coefficients
at every cusp, with cusp-width factors.
The level-one formula \(k=f(0,0)/2\) therefore becomes
\[
 k=\frac12\sum_{\mathfrak c}
             \frac{h_{\mathfrak c}}{N_{\mathfrak c}}\,
                         c_{\mathfrak c}(0,0),
\]
in the notation of the cited product theorem.
Here \(h_{\mathfrak c}\) is the cusp width and
\(N_{\mathfrak c}=N/e\) for a cusp represented by \(f/e\).
All cusp data belong to the arithmetic input.

Integral product exponents define a minimal super vector
space and its formal Fock character in each row.
Under an absolute coefficient bound they also define
the trace-class diagonal construction used above.
A reflective denominator presentation supplies further
Lie and parity data. Its ordinary quiver Hall algebra
retains the even Cartan construction, while the super
relations remain separately specified.
Neither the eight scalar identities nor a common Cartan
matrix identifies those two algebras.

\section{The compact comparison}

The compact constructions now have explicit coefficients,
units and normalizations. The root determinant uses signed
root multiplicities and the chosen Weyl monomial.
The Hilbert-scheme state inserts the tangent characteristic
class. The common-torsor extension product uses flags over
\(BE\), and its critical coefficients require orientation
and base-change maps on those flags.

The remaining compact chiral comparison must act on state
complexes defined independently from the Hall image.
It must construct diagonal and Gysin maps preserving normal
oscillators, cycle joining, charge, vacuum insertions and
higher coordinate gluing.
At integral specialization it must account for the
\(t^4\) and \(t^2\) torsion of
Proposition~\ref{mht:prop:compact-smith}.
An operator comparison must also retain the determinant
cross lines and the distinction between local and
elliptically wrapped support.
The scalar Igusa identities provide exact tests for such
a map. They do not determine the map.
